\documentclass[11pt]{article}
\usepackage[T1]{fontenc}
\usepackage{lmodern,microtype,amsmath,amssymb,amsthm,mathtools}
\usepackage[hmargin=1.04in,vmargin=0.92in]{geometry}
\usepackage[hidelinks]{hyperref}
\setlength{\emergencystretch}{1.3em}
\hypersetup{pdftitle={How closely can the zeros of two Dirichlet L-functions be matched?}}
\newtheorem{theorem}{Theorem}[section]
\newtheorem{lemma}[theorem]{Lemma}
\newtheorem{corollary}[theorem]{Corollary}
\newcommand{\E}{\mathbb E}
\newcommand{\norm}[1]{\lVert#1\rVert}
\newcommand{\dd}{\,\mathrm d}
\title{How closely can the zeros of two Dirichlet\\ \(L\)-functions be matched?}
\author{Research note}
\date{27 November 2026}
\begin{document}
\maketitle
\begin{abstract}
Count nontrivial zero ordinates of two distinct primitive Dirichlet characters
in \((T,2T]\) with multiplicity. We prove that the proportion left unmatched by the
best one-to-one matching of displacement at most \(K/\log T\) tends to zero
if and only if \(K\to\infty\). Even an allowance of one hundred mean spacings
leaves a positive proportion unmatched on each side. The dichotomy holds for
different finite products of the same degree, with arbitrary overlap among their
primitive factors. The upper limiting proportion decays faster than every inverse power of \(K\). The lower limiting proportion decays more slowly than every \(e^{-sK^2}\), \(s>0\). We also show that
binning at bounded microscopic resolution retains Hellinger information of order
\(T\log T\). The arithmetic input is the unconditional short-interval
joint central limit theorem of Bombieri and Perelli. Dirichlet GRH identifies
distances of ordinates with distances of the zeros themselves.
\end{abstract}

\section{The matching threshold}

Let \(\chi_1,\ldots,\chi_d\) be distinct fixed primitive Dirichlet characters,
possibly including the conductor-one character. For \(a\in\mathbb Z_{\ge0}^d\)
let \(Z_a(T)\) be the multiset of ordinates in \((T,2T]\) of the nontrivial zeros
of \(\prod_{j=1}^d L(s,\chi_j)^{a_j}\). A zero of order \(m\) in the \(j\)th
factor contributes \(a_jm\) occurrences. A partial matching pairs occurrences of
\(Z_a(T)\) with occurrences of \(Z_b(T)\), using each occurrence at most
once. Write \(D_T(K)\) for the minimum total number unpaired, summed over both
sides, among matchings satisfying \(|\gamma-\gamma'|\le K/\log T\) on every
pair. The combined list of two degree-one factors has size
\((1/\pi+o(1))T\log T\). Within one such list the mean spacing is asymptotic
to \(2\pi/\log T\), so one hundred mean spacings corresponds to \(K=200\pi\).

\begin{theorem}\label{th:threshold}
Suppose \(a\ne b\) and \(\sum a_j=\sum b_j=r\). For every fixed finite \(K\ge0\),
\[
 \liminf_{T\to\infty}\frac{D_T(K)}{T\log T}>0 .
\]
For every function \(K(T)\ge0\),
\[
 D_T(K(T))=o(T\log T)\qquad\Longleftrightarrow\qquad K(T)\longrightarrow\infty .
\]
Every fixed-allowance matching leaves a positive proportion of each list
unpaired. For a diverging allowance an order-preserving partial matching
attains the conclusion.
\end{theorem}

If the degrees differ, the absolute difference of the list sizes is
\((|\sum a_j-\sum b_j|/(2\pi)+o(1))T\log T\), a lower bound at every
allowance. If \(a=b\), the identity gives \(D_T(0)=0\). Distinct conjugate
characters are allowed. The lists use positive ordinates and preserve the
distinction lost by folding the two half-planes together. Shared primitive
factors may participate in the matching in any way.

Fujii compared exact coincidences, gaps and rank discrepancies
for primitive characters \cite{fujii}. Bombieri and Perelli obtained a
positive proportion of multiplicity differences for a wider class of
\(L\)-functions \cite[Theorem 1]{bp}. ProveIt's article on Dirichlet GRH
uses the joint short-interval law to show that the zero multiplicity vectors have
full rank even after a positive proportion of the observations are removed.
It also gives a positive cost for moving or deleting zeros in intervals of
triple-logarithmic length \cite{repo}. Theorem~\ref{th:threshold} allows
reassignment of the entire dyadic list and determines when the discarded
proportion vanishes. Several nearby occurrences may compete for one partner,
so the discard count has a different meaning from nearest-neighbour distance.

\section{Short-interval zero counts}

Write \(N_j(t)\) for the number of nontrivial zeros at positive ordinates up to
\(t\), and \(C_j(t,h)=N_j(t+h)-N_j(t)\). Put \(C_a=\sum_j a_j C_j\)
and \(c=a-b\), so \(\sum_j c_j=0\).
Here \(t\) is uniform on \([T,2T]\). Values at individual endpoints have no
effect on these integrals.

\begin{lemma}\label{lem:CLT}
Let \(M=M(T)\to\infty\), \(M\le \log T/\log\log T\), and \(h=M/\log T\).
Then
\begin{align}
 \frac{C_a(t,h)-C_b(t,h)}{\sqrt{\log M}}
 &\ \Longrightarrow\ G_c, &
 G_c&\sim N(0,\norm c_2^2/\pi^2),\label{eq:clt}\\
 \E|C_a(t,h)-C_b(t,h)|
 &= (\eta_c+o(1))\sqrt{\log M},&
 \eta_c&=\frac{\sqrt2\,\norm c_2}{\pi^{3/2}}.\label{eq:L1}
\end{align}
For \(v=a,b\), we also have \(C_v(t,h)/M\to r/(2\pi)\) in probability.
\end{lemma}

\begin{proof}
Write \(S_j(t)=\pi^{-1}\arg L(1/2+it,\chi_j)\) with the standard continuation
of the argument. Stirling's formula in the Riemann--von Mangoldt formula gives,
for almost every \(t\in[T,2T]\) and \(0<h\le1\),
\[
 C_j(t,h)=\frac{h}{2\pi}\log T+S_j(t+h)-S_j(t)+O_{\chi_j}(h+1/T).
\]
The conductor term and \(h\log(t/T)\) are included in the error.
Also \(N_j(2T)-N_j(T)=(T/(2\pi))\log T+O_{\chi_j}(T)\), and
\(N_j(u+1)-N_j(u)=O_{\chi_j}(\log T)\) at heights \(u\asymp T\).

Theorem~2 of Bombieri and Perelli \cite{bp} applies to these factors.
Their Euler parameters have modulus at most one, and they have the classical
functional equations and continuations. The prime number theorem for fixed
progressions modulo the least common multiple of the conductors gives
Hypothesis B(IV) with diagonal parameter one, by character orthogonality.
Their Hypothesis D is the classical density estimate for these functions
cited on p.~273 of the same paper. Dirichlet GRH also implies this hypothesis.
The imaginary parts of their normalized log increments \(V_j\) converge jointly
to independent normals of variance \(1/(2\pi)\). Their normalization
\(\Delta_h S_j=\sqrt{2\pi\log M}\,\operatorname{Im}V_j/\pi\)
gives \eqref{eq:clt} and the stated convergence in probability.

For the first absolute moment, p.~279 of \cite{bp} approximates \(V_j\) by
normalized Dirichlet polynomials \(U_j\) satisfying
\(\E|V_j-U_j|\ll(\log M)^{-1/4}\).
Their Lemma~2 bounds \(\E|U_j|^2\) uniformly. The same second-moment
bound holds for every fixed linear combination of their imaginary parts.
Bounded truncation and then uniform integrability give convergence of its
first absolute moment. The displayed \(L^1\) approximation transfers this
conclusion to the log increments. Finally \(\E|G_c|=\eta_c\).
\end{proof}

\section{A fixed displacement leaves a positive proportion}

Let \(U_A,U_B\) count the discarded occurrences of a matching, and \(P\)
its pairs. For \(t\in[T,2T-h]\), put \(W_t=(t,t+h]\).
A point \(x\) belongs to \(W_t\) for starts in \([x-h,x)\).
Two points at distance \(d\) have different membership for a set of starts of
measure \(2\min(d,h)\) on the whole real line. Integrating over the smaller
range of starts therefore gives
\begin{equation}\label{eq:transport}
 \int_T^{2T-h}|C_a(t,h)-C_b(t,h)|\dd t
 \le h(U_A+U_B)+2\sum_{\rm pairs}\min(h,|\gamma-\gamma'|).
\end{equation}
Indeed each discarded occurrence contributes for at most \(h\) units of
time, and the triangle inequality sums the pairwise differences.
Omitting the last interval of starts from \eqref{eq:L1} costs
\(O(h\log T)\) in the unnormalized integral by the unit-interval bound.

If the first assertion of Theorem~\ref{th:threshold} failed, there would be a
sequence of heights along which \(d_T=D_T(K)/(T\log T)\to0\).
Choose
\[
 M=\min\left\{\frac{\log T}{\log\log T},\,
                  (d_T+1/(T\log T))^{-1/2}\right\}.
\]
Then \(M\to\infty\) and \(Md_T\to0\). Use an optimal matching in
\eqref{eq:transport}, divide by \(T\), and use \(P=O_{a,b}(T\log T)\).
The right side is at most \(Md_T+O_{a,b}(K)\), while the left side is
asymptotic to \(\eta_c\sqrt{\log M}\to\infty\). Moreover
\(U_A-U_B=|Z_a(T)|-|Z_b(T)|=O_{a,b}(T)\), so each side supplies a positive
proportion. A bounded subsequence of any allowance also precludes convergence
of the discarded proportion to zero, by monotonicity.

\section{A construction at every growing allowance}

\begin{theorem}\label{th:upper}
If \(K=K(T)\to\infty\), put
\(m=\min\{K,\log T/\log\log T\}\). Then
\[
 \frac{D_T(K)}{T\log T}
 \le (\eta_c+o(1))\frac{\sqrt{\log m}}{m}.
\]
Consequently \(g(K)=\limsup_{T\to\infty}D_T(K)/(T\log T)\) satisfies
\(\limsup_{K\to\infty}K g(K)/\sqrt{\log K}\le\eta_c\).
\end{theorem}

\begin{proof}
Use cells of width \(h=m/\log T\) with origin \(T+u\), where
\(u\) is uniform on \([0,h)\). Within each complete cell match in increasing
order as many points as possible. The unpaired count there is
\(|C_a(t,h)-C_b(t,h)|\) at its left endpoint \(t\).
The two partial boundary cells contain \(O(\log T)\) occurrences in total and
may be discarded. Averaging the sum over complete cells gives exactly
\[
 \frac1h\int_T^{2T-h}|C_a(t,h)-C_b(t,h)|\dd t .
\]
Some origin attains at most this average. Apply \eqref{eq:L1}. The boundary
term, divided by \(T\log T\), is negligible throughout the stated range. This
proves the bound and completes Theorem~\ref{th:threshold}.
For the iterated limit, any violating sequence \(K_n\to\infty\) could be
tested at arbitrarily large heights \(T_n\) with
\(K_n\le\log T_n/\log\log T_n\), approaching the corresponding upper limits
to within \(o(\sqrt{\log K_n}/K_n)\). The first bound rules this out.
\end{proof}

An exact linear-time algorithm is available for sorted lists. Compare their
two smallest remaining entries. Discard the smaller one if the gap exceeds
the allowance, and otherwise match the pair. In the first case the smaller
entry has no possible partner. In the second case an optimal matching can use
the proposed pair. If exactly one is matched, replace its partner.
If both have different partners \(a'\ge a\) and \(b'\ge b\), replace the
two pairs by \((a,b),(a',b')\). From \(|a-b'|,|a'-b|\le\delta\) and the order,
\(|a'-b'|\le\delta\). If both were unmatched, adding them would improve the
matching. Induction proves optimality, including ties.
The accompanying \texttt{check\_matching.py} tests this algorithm against
unrestricted bipartite matching and verifies both integrated window inequalities
by exact integer arithmetic.

\section{Decay faster than any inverse power}

The one-cell construction gives a uniform bound throughout a growing
mesoscopic range. Using many cells at once yields a stronger estimate
after the height tends to infinity at each fixed allowance.

\begin{theorem}\label{th:rapid}
For \(g(K)=\limsup_{T\to\infty}D_T(K)/(T\log T)\) and every \(A>0\),
\[
                    g(K)\ll_{A,a,b} K^{-A}\qquad(K\ge2).
\]
\end{theorem}

\begin{proof}
Fujii's Main Theorem \cite[pp.~225--226]{fujiiI}, with its corrections
\cite{fujiiII}, gives the following fixed-moment bound.
For fixed primitive \(\chi\), fixed integer \(k\ge1\), dyadic averaging
interval of length \(T\), and \(0<u\le 1/\log\log T\),
\[
 \frac1T\int_T^{2T}|S_\chi(t+u)-S_\chi(t)|^{2k}\dd t
       \ll_{k,\chi} \{\log(2+u\log T)\}^{k}.
\]
Only bounded shifts and fixed conductor and moment order occur here.
The zero-count formula and the triangle inequality in \(L^{2k}\) yield
for \(v=a,b\)
\begin{equation}\label{eq:highmoment}
 \E|C_v(t,u)-ru\log T/(2\pi)|^{2k}
 \ll_{k,a,b}\{\log(2+u\log T)\}^{k}.
\end{equation}

Fix \(B>A\). Let \(J=\lceil K^B\rceil\), \(h=K/(4\log T)\),
\(H=Jh\), and \(\mu=rK/(8\pi)\).
At each fixed \(K\), take \(T\) sufficiently large that \(H\le1/\log\log T\).
Divide a block \((t,t+H]\) into \(J\) cells of width \(h\).
Call the block good if, for \(v=a,b\) and \(1\le j\le J\),
\[
 |C_v(t,jh)-j\mu|\le\mu/4 .
\]
For a good block, pair points in increasing rank through the end of the
shorter list. A point in cell \(j\) has rank at most \((j+1/4)\mu\)
and greater than \((j-1-1/4)\mu\) if \(j>1\).
The other list's count through cell \(j+1\), when present, is at least
\((j+1-1/4)\mu\). Its count through cell \(j-2\), when present, is at most
\((j-2+1/4)\mu\). Thus matching equal ranks places them in the same
or adjacent cells, with displacement at most \(2h\). The two totals differ
by at most \(\mu/2\), which bounds the number discarded.

For \(t\) uniform on \([T,2T]\), the union bound and \eqref{eq:highmoment}
show that the probability \(p\) of a bad block satisfies
\[
 p\ll_{k,a,b} J\{\log(2+JK)\}^{k}/K^{2k}.
\]
Discard all occurrences in a bad block. By \eqref{eq:highmoment} at \(k=1\)
the square mean of the total block count is \(O_{a,b}((JK)^2)\).
Cauchy--Schwarz bounds its contribution on the bad event by
\(O_{a,b}(JK\sqrt p)\) in expectation.

Tile by blocks of width \(H\) with a uniform random grid origin and
discard the two partial boundary blocks. They contain \(O_{a,b}(\log T)\)
points because \(H\le1\). As in Theorem~\ref{th:upper}, the expected cost
on complete blocks is bounded by \(T/H\) times the preceding upper bound
on the expected discard count of a single block. Consequently
\[
 g(K)\ll_{k,a,b} J^{-1}
       +J^{1/2}\{\log(2+JK)\}^{k/2}K^{-k}.
\]
Take an integer \(k>A+B/2+1\) to obtain the claim.
\end{proof}

\section{A lower bound as the allowance grows}

Write \(x_+=\max(x,0)\) and \(f(K)=\liminf_{T\to\infty}D_T(K)/(T\log T)\).
\begin{theorem}\label{th:tails}
For every \(s>0\),
\begin{equation}\label{eq:gausslower}
 \liminf_{K\to\infty}\frac{e^{sK^2}}K f(K)
 \ \ge\ \E\bigl[(\sqrt{s}|G_c|-r/\pi)_+\bigr]>0 .
\end{equation}
In particular \(\lim_{K\to\infty}K^{-2}\log f(K)=0\).
\end{theorem}

\begin{proof}
Fix a sequence \(K_n\to\infty\). Choose heights \(T_n\to\infty\) so
that \(M=e^{sK_n^2}\le\log T_n/\log\log T_n\) and
\[
 D_{T_n}(K_n)/(T_n\log T_n)\le f(K_n)+o(K_n/M).
\]
This is possible by definition of the lower limit. Write \(K=K_n\),
\(T=T_n\), \(h=M/\log T\), \(\delta=K/\log T\), and take starts uniform
on \(J=[T+\delta,2T-h-\delta]\). Restricting the uniform dyadic law to \(J\)
changes probabilities by \(O((h+2\delta)/T)\).

Put \(W=(t,t+h]\) and \(W^+=(t-\delta,t+h+\delta]\).
Every possible partner of an occurrence in \(W\) lies in \(W^+\).
Let \(B_\partial\) be the total number of \(b\)-occurrences in the two
pieces of \(W^+\setminus W\) and define \(A_\partial\) similarly. If
\(X=C_a(t,h)-C_b(t,h)\), the number of discarded occurrences within \(W\)
is at least
\[
 (X-B_\partial)_+ +(-X-A_\partial)_+ .
\]
Every discarded occurrence contributes for a set of starts of measure at
most \(h\). Multiplying the expectation of this bound by
\(|J|/h\) thus gives a lower bound for \(D_T(K)\).

By \eqref{eq:clt}, \(X/K\Longrightarrow\sqrt{s}G_c\). Each flank count
of either product, divided by \(K\), tends in probability to \(r/(2\pi)\)
by Lemma~\ref{lem:CLT} at parameter \(K\). Its start is the uniform
distribution on \(J\) translated by \(-\delta\) or \(h\), at total variation
distance \(o(1)\) from the uniform dyadic law.
Hence \(A_\partial/K,B_\partial/K\to r/\pi\) in probability.
Slutsky's theorem gives joint convergence with \(X/K\).
Taking expectations of bounded truncations of the positive parts and then
increasing the truncation yields
\[
 \liminf_{n\to\infty}\frac{M}{K}\frac{D_T(K)}{T\log T}
 \ge \E\bigl[(\sqrt{s}|G_c|-r/\pi)_+\bigr].
\]
The expectation is positive since \(G_c\) has nonzero variance and unbounded
support. Our choice of \(T_n\) proves \eqref{eq:gausslower}. It gives a
lower limit of at least \(-s\) for \(\log f(K)/K^2\) for every \(s>0\).
The upper limit is at most zero since \(f\) is bounded.
\end{proof}

\section{Information after binning}

For the oriented Li zero-angle process in ProveIt, the number of jumps at an
atom has a Poisson law whose mean is observation time times multiplicity.
The angle on one oriented branch determines the ordinate. Binning these
angles amounts to observing count histograms of the ordinate band.

\begin{corollary}\label{cor:hell}
Partition \((T,2T]\) into disjoint measurable bins of diameter at most
\(K/\log T\) for fixed \(K\). Write \(A_i,B_i\) for their respective counts.
Then, uniformly over such partitions,
\[
 \sum_i(\sqrt{A_i}-\sqrt{B_i})^2\asymp_{a,b,K}T\log T .
\]
For a diverging allowance there are interval partitions respecting this
diameter bound for which the same sum is \(o(T\log T)\).
\end{corollary}

\begin{proof}
Pairing entries within each bin gives \(\sum_i|A_i-B_i|\ge D_T(K)\).
Cauchy--Schwarz gives
\[
 \frac{(\sum_i|A_i-B_i|)^2}{2\sum_i(A_i+B_i)}
 \le\sum_i(\sqrt{A_i}-\sqrt{B_i})^2\le\sum_i|A_i-B_i|.
\]
The fixed-allowance conclusion follows from Theorem~\ref{th:threshold} and
the total count. The partition of Theorem~\ref{th:upper} gives the other
conclusion, including its two boundary cells.
\end{proof}

For the article's symmetric jump process, the rate at a positive angle
\(\theta=2\arctan(1/(2\gamma))\) is
\((m_\chi(\gamma)+m_{\bar\chi}(\gamma))/2\), by conjugation of the zeros.
Take \(a=e_\chi+e_{\bar\chi}\) and \(b=e_\psi+e_{\bar\psi}\), allowing repeated
indices for a real character. These are distinct equal-degree products exactly
when the two conjugacy orbits differ. Halving both histograms halves their
squared-root discrepancy. Counting both angle signs restores the same factor
of two. Consequently the full-band Hellinger sum in the article has order
\(T\log T\) for all nonconjugate pairs, even after binning at fixed microscopic
resolution in the ordinate. The article's single-prime estimate gives
a lower bound of order \(T/\log T\); this argument gains two powers of
\(\log T\). On the band, \(|\theta'(\gamma)|=1/(\gamma^2+1/4)\asymp T^{-2}\),
so angle bins of width \(O(T^{-2}/\log T)\) suffice.

For independent Poisson counts of means \(uA_i\) or \(uB_i\), their joint
Hellinger affinity is exactly \(\exp[-(u/2)\sum_i(\sqrt{A_i}-\sqrt{B_i})^2]\),
by multiplicativity. Thus bounded microscopic bin widths preserve the order
of this exponent for every choice of bin edges.

\paragraph{Further questions.}
Does \(D_T(K)/(T\log T)\) have a limit for each fixed \(K>0\)? How does it
decay as \(K\) grows? The interval between Theorems~\ref{th:rapid}
and~\ref{th:tails} leaves room for different rates. A sharper answer would
quantify how often a local excess of one spectrum can be absorbed by unused
partners nearby. The density of common zeros of two distinct primitive
\(L\)-functions remains open in general.

\begin{thebibliography}{9}
\bibitem{fujii} A. Fujii, \emph{On the zeros of Dirichlet \(L\)-functions (V)},
Acta Arith. \textbf{28} (1976), 395--403.
\href{https://doi.org/10.4064/aa-28-4-395-403}{doi:10.4064/aa-28-4-395-403}.
\bibitem{fujiiI} A. Fujii, \emph{On the zeros of Dirichlet \(L\)-functions. I},
Trans. Amer. Math. Soc. \textbf{196} (1974), 225--235.
\href{https://doi.org/10.1090/S0002-9947-1974-0349603-2}{doi:10.1090/S0002-9947-1974-0349603-2}.
\bibitem{fujiiII} A. Fujii, \emph{On the zeros of Dirichlet \(L\)-functions. II
(with corrections to Part I and the subsequent papers)},
Trans. Amer. Math. Soc. \textbf{267} (1981), 33--40.
\href{https://doi.org/10.1090/S0002-9947-1981-0621970-5}{doi:10.1090/S0002-9947-1981-0621970-5}.
\bibitem{bp} E. Bombieri and A. Perelli, \emph{Distinct zeros of \(L\)-functions},
Acta Arith. \textbf{83} (1998), 271--281.
\href{https://doi.org/10.4064/aa-83-3-271-281}{doi:10.4064/aa-83-3-271-281}.
\bibitem{repo} V. Reshetnikov, \emph{ProveIt}, Analysis/GeneralizedRiemannHypothesis,
snapshot \texttt{cc919c585f90}. See
\href{https://github.com/VladimirReshetnikov/ProveIt/blob/cc919c585f90be8e0fba7a1b90c70afdea39998e/Analysis/GeneralizedRiemannHypothesis/Article/sections/rec.tex}{\texttt{rec.tex}, A cost for moving or deleting zeros, and A positive proportion of distinct pole ordinates}, and
\href{https://github.com/VladimirReshetnikov/ProveIt/blob/cc919c585f90be8e0fba7a1b90c70afdea39998e/Analysis/GeneralizedRiemannHypothesis/Article/sections/li.tex}{\texttt{li.tex}, Finite-band information and testing}.
\end{thebibliography}
\end{document}
